\begin{objectivebox}[이 장에서 붙잡을 질문]
사람에게는 당연해 보이는 계산 하나를 컴퓨터가 수행하려면, 무엇을 결정해야 하며 그 결정에는 어떤 비용이 따를까?
\end{objectivebox}

\section*{1. 한 줄은 목적지를 말한다}

이 책은 다음과 같은 작은 계산에서 출발한다.

\[
  A=\begin{bmatrix}1&2\\3&4\end{bmatrix},
  \qquad
  B=\begin{bmatrix}5&6\\7&8\end{bmatrix}
\]

\[
  C=A\times B=
  \begin{bmatrix}19&22\\43&50\end{bmatrix}.
\]

두 숫자 표 \(A\)와 \(B\)를 곱해 \(C\)를 얻었다. 예를 들어 \(C\)의 왼쪽 위 칸 \(19\)는 \(A\)의 첫 가로줄 \((1,2)\)와 \(B\)의 첫 세로줄 \((5,7)\)에서 \(1\times5+2\times7=19\)로 계산한다. 나머지 세 칸도 같은 규칙으로 얻는다.

이 식은 어떤 답을 얻어야 하는지를 정한다. 그러나 그 답을 어떻게 계산할지까지 정해 주지는 않는다. 작은 계산을 몇 개 만들지, 어느 계산을 먼저 할지, 중간값을 어디에 둘지는 별도로 결정해야 한다. 같은 목적지로 가는 길이 하나뿐이라고 쓰여 있지 않은 셈이다.

\begin{bookfigure}{한 줄과 유한한 기계}{fig:preface-equation-and-machine}
\begin{tikzpicture}[diagram]
  \path[use as bounding box] (0,0) rectangle (13.0,8.45);

  \draw[diagram panel] (0.05,0.45) rectangle (4.10,8.30);
  \node[diagram panel title] at (0.30,7.96) {수학이 정한 것};
  \node[font=\sffamily\bfseries\fontsize{19pt}{21pt}\selectfont] at (2.08,5.35)
    {\(C=A\times B\)};
  \node[diagram note,align=center] at (2.08,4.25)
    {두 숫자 표를 정해진 규칙으로 결합해\\결과 (C)를 만든다};
  \node[concept emphasis,minimum width=32mm,minimum height=11mm] at (2.08,2.15)
    {원하는 결과의 규칙};
  \node[diagram note,align=center,text=black!65] at (2.08,1.15)
    {배치·순서·저장 위치는\\아직 정하지 않았다};

  \draw[concept dashed arrow] (4.27,5.95) -- (5.12,5.95);
  \draw[concept dashed arrow] (4.27,2.75) -- (5.12,2.75);
  \node[diagram note,align=center] at (4.70,4.36)
    {한 가지\\계획으로\\결정되지 않음};

  \draw[diagram panel,densely dashed] (5.30,0.45) rectangle (12.95,8.30);
  \node[diagram panel title] at (5.55,7.96) {기계에서 따로 정할 것};

  \node[concept emphasis,minimum width=28mm,minimum height=13mm] (ops) at (7.12,6.35)
    {연산기 수\\동시에 몇 개?};
  \node[concept box,minimum width=28mm,minimum height=13mm] (time) at (10.92,6.35)
    {신호·연산 시간\\0이 아님};
  \node[concept box,minimum width=28mm,minimum height=13mm] (path) at (7.12,4.33)
    {이동 거리·통로 폭\\한 번에 얼마나?};
  \node[concept box,minimum width=28mm,minimum height=13mm] (store) at (10.92,4.33)
    {가까운 저장공간\\얼마나 둘까?};
  \node[concept box,minimum width=28mm,minimum height=13mm] (bits) at (7.12,2.30)
    {값의 비트 수\\범위와 정밀도};
  \node[concept emphasis,minimum width=28mm,minimum height=13mm] (plan) at (10.92,2.30)
    {실행 계획\\배치·순서·재사용};

  \draw[concept arrow] (ops.east) -- (time.west);
  \draw[concept arrow] (path.east) -- (store.west);
  \draw[concept arrow] (bits.east) -- (plan.west);

  \node[diagram footer] at (9.12,0.83)
    {유한한 자원을 어떻게 엮느냐에 따라 같은 결과의 비용이 달라진다};
\end{tikzpicture}
\end{bookfigure}


\figref{fig:preface-equation-and-machine}의 왼쪽은 수학이 고정한 규칙이고, 오른쪽은 기계에서 별도로 정해야 할 조건이다. 가운데의 점선은 식에서 기계로 값이 흘러간다는 배선이 아니다. \emph{이 한 줄만으로 오른쪽의 배치가 한 가지로 결정되지 않는다}는 경계를 나타낸다.

\section*{2. 유한한 기계가 덧붙이는 다섯 조건}

종이 위의 식에는 연산기가 몇 개인지 적혀 있지 않다. 신호가 이동하는 데 걸리는 시간도, 값을 가까이 보관할 공간도 보이지 않는다. 실제 기계는 이 빈칸을 유한한 자원으로 채운다.

첫째, 동시에 쓸 수 있는 \keyterm{연산} 장치의 수가 유한하다. 해야 할 작은 계산이 장치보다 많다면 일부는 차례를 기다려야 한다. 둘째, 신호 전달과 연산에는 0이 아닌 시간이 든다. 앞선 결과가 있어야 시작할 수 있는 계산에는 \keyterm{의존성}이 생긴다. 셋째, 값을 옮길 거리와 통로의 폭이 유한하다. 연산이 빨라도 필요한 값이 제때 도착하지 않으면 장치는 기다린다.

넷째, 연산기 가까이에 둘 수 있는 저장공간이 유한하다. 자주 쓰는 값이 멀리 있다면 \keyterm{데이터 이동}이 늘어난다. 다섯째, 값 하나에 쓸 수 있는 비트 수가 유한하다. 이 제한은 표현할 수 있는 범위와 \keyterm{정밀도}를 정하고, 때로는 계산 순서에 따라 기계가 내는 값에도 영향을 준다.

이 다섯 조건을 붙이면 한 줄의 식이 답하지 않던 질문이 드러난다.

\begin{center}
\small
\begin{tabular}{>{\raggedright\arraybackslash}p{0.42\textwidth} >{\raggedright\arraybackslash}p{0.39\textwidth}}
\toprule
\sffamily\bfseries 수학식만으로 정해지지 않는 질문 & \sffamily\bfseries 기계에서 확인할 조건 \\
\midrule
작은 계산으로 어떻게 나눌 것인가 & 같은 함수를 분해하는 방법과 조합 비용 \\
무엇이 앞선 결과를 기다려야 하는가 & 0이 아닌 신호·연산 시간 \\
몇 개를 동시에 계산할 것인가 & 동시에 쓸 수 있는 연산기 수 \\
값을 어디에 두고 얼마나 옮길 것인가 & 이동 거리·통로 폭·가까운 저장 용량 \\
한 값을 몇 비트로 나타낼 것인가 & 값 하나에 허용한 비트 수 \\
\bottomrule
\end{tabular}
\end{center}

첫 질문은 기계가 작아서만 생기는 문제가 아니다. 같은 함수를 작은 함수로 나누는 방법 자체가 여러 가지일 수 있다. 나머지 네 질문은 그 여러 방법을 유한한 기계에 놓을 때 서로 다른 비용을 만든다.

\section*{3. 계산의 요구량과 기계의 자원은 다르다}

앞으로는 계산과 기계를 한 덩어리로 세지 않는다. 먼저 계산이 무엇을 요구하는지 적고, 그다음 기계가 무엇을 제공하는지 적는다.

\begin{center}
\small
\begin{tabular}{>{\raggedright\arraybackslash}p{0.40\textwidth} >{\raggedright\arraybackslash}p{0.41\textwidth}}
\toprule
\sffamily\bfseries 계산이 요구하는 양 & \sffamily\bfseries 기계가 제공하는 자원 \\
\midrule
해야 할 연산의 수 & 동시에 사용할 수 있는 연산기 \\
반드시 이어지는 의존 단계 & 실행에 쓸 수 있는 시간 \\
옮겨야 할 데이터의 양 & 통로의 폭과 이동 거리 \\
살려 두어야 할 중간값 & 가까운 저장공간의 용량 \\
결과에 필요한 표현 범위와 정확도 & 값 하나에 배정한 비트 수 \\
\bottomrule
\end{tabular}
\end{center}

예를 들어 해야 할 연산이 백 개라는 사실과 연산기가 백 개라는 사실은 전혀 다른 문장이다. 연산 백 개를 연산기 하나로 차례로 실행할 수도 있고, 여러 연산기에 나누어 동시에 실행할 수도 있다. 두 방법이 같은 답을 낼 수는 있지만 걸리는 시간, 필요한 회로, 값을 옮기는 양은 같지 않다.

따라서 이 책에서 `빠르다'라는 말은 혼자 서지 않는다. 무엇을 더 썼는지 함께 묻는다. 더 많은 연산기를 썼는가, 더 넓은 통로를 썼는가, 더 가까운 곳에 값을 복사했는가, 아니면 정밀도나 프로그램 가능성을 줄였는가. 속도·면적·전력·제어의 복잡성을 하나의 막연한 점수로 합치지 않고 각각 따로 살펴볼 것이다.

\begin{shadequote}
수학은 원하는 결과의 규칙을 정하지만, 유한한 기계에서 그 결과를 얻는 비용을 한 가지로 고정하지는 않는다.
\end{shadequote}

\section*{4. 내려가서 이해하고, 다시 올라가서 배치한다}

그렇다면 어디서부터 시작해야 할까. 복잡한 계산을 한꺼번에 기계에 얹으려 하면 식과 장치 사이가 너무 멀다. 이 책은 먼저 아래로 내려가 가장 작은 규칙을 찾고, 그 규칙을 다시 조합하면서 위로 올라온다.

\begin{bookfigure}{내려갔다 올라오는 두 길}{fig:preface-down-and-up}
\begin{tikzpicture}[diagram]
  \path[use as bounding box] (0,0) rectangle (13.0,11.75);

  \draw[diagram panel] (0.05,0.35) rectangle (4.25,11.60);
  \node[diagram panel title] at (0.30,11.26) {내려가는 길};
  \node[diagram note,anchor=west,text=black!65] at (0.30,10.80) {작은 규칙을 찾는다};

  \node[concept emphasis,minimum width=31mm] (down1) at (2.15,9.63)
    {행렬 결과의 한 부분};
  \node[concept box,minimum width=31mm] (down2) at (2.15,7.62)
    {더 작은 산술};
  \node[concept box,minimum width=31mm] (down3) at (2.15,5.61)
    {이진 덧셈};
  \node[concept emphasis,minimum width=31mm,minimum height=11mm] (down4) at (2.15,3.45)
    {진리표\\제1장의 출발점};
  \draw[concept arrow] (down1.south) -- (down2.north);
  \draw[concept arrow] (down2.south) -- (down3.north);
  \draw[concept arrow] (down3.south) -- (down4.north);
  \node[diagram note,align=center,text=black!65] at (2.15,1.55)
    {복잡한 계산을\\가장 작은 명세까지 분해};

  \draw[diagram panel] (4.55,0.35) rectangle (12.95,11.60);
  \node[diagram panel title] at (4.80,11.26) {올라오는 길};
  \node[diagram note,anchor=west,text=black!65] at (4.80,10.80) {유한한 자원에 배치한다};

  \node[concept box,minimum width=64mm,minimum height=10mm] (up6) at (8.75,9.72)
    {6 · 전용화 \quad 행렬 경로·AI 가속기 \quad 18--20장};
  \node[concept box,minimum width=64mm,minimum height=10mm] (up5) at (8.75,8.12)
    {5 · 데이터 재사용 \quad 메모리·타일링 \quad 13--17장};
  \node[concept emphasis,minimum width=64mm,minimum height=10mm] (up4) at (8.75,6.52)
    {4 · 복제·공유 \quad 병렬 실행·GPU \quad 8--12장};
  \node[concept box,minimum width=64mm,minimum height=10mm] (up3) at (8.75,4.92)
    {3 · 시간적 재사용 \quad 상태·명령 \quad 6--7장};
  \node[concept box,minimum width=64mm,minimum height=10mm] (up2) at (8.75,3.32)
    {2 · 조합 \quad 덧셈·곱셈·ALU \quad 2--5장};
  \node[concept emphasis,minimum width=64mm,minimum height=10mm] (up1) at (8.75,1.72)
    {1 · 명세 \quad 진리표 \quad 1장};

  \draw[concept arrow] (up1.north) -- (up2.south);
  \draw[concept arrow] (up2.north) -- (up3.south);
  \draw[concept arrow] (up3.north) -- (up4.south);
  \draw[concept arrow] (up4.north) -- (up5.south);
  \draw[concept arrow] (up5.north) -- (up6.south);

  \draw[concept dashed arrow] (down4.east) -- ++(0.34,0) |- (up1.west);
  \node[diagram note,anchor=west,text=black!65] at (4.78,0.72)
    {역사 연표가 아니라, 비용을 비교하기 위한 설명 순서};
\end{tikzpicture}
\end{bookfigure}


\figref{fig:preface-down-and-up}의 왼쪽 길은 계산을 이해하기 위한 분해다. 행렬 결과의 한 부분을 더 작은 산술로, 산술을 이진 덧셈으로, 덧셈을 참과 거짓의 표로 내려간다. 이 길의 끝인 \keyterm{진리표}가 제1장의 실제 출발점이다.

오른쪽 길은 작은 규칙을 유한한 자원에 놓는 과정이다. 진리표로 기능을 명세하고, 작은 기능을 덧셈·곱셈·산술논리장치로 조합한다. 이어 상태와 명령을 도입하면 같은 회로를 시간에 따라 다시 쓸 수 있다. 이 다리가 있어야 회로를 무작정 복제하는 방법과 반복해서 사용하는 방법을 비교할 수 있다.

그 위에서 계산을 복제하고 자원을 공유하면 병렬 실행과 GPU가 나타난다. 하지만 연산기 수만 늘린다고 끝나지 않는다. 필요한 값을 가까이 두고 되쓰는 메모리 구조와 타일링을 거쳐야 한다. 마지막에는 프로그램 가능성을 일부 줄이고 특정 계산 경로를 강화한 행렬 연산 장치와 AI 가속기를 살펴본다. 이 순서는 제품의 역사 연표가 아니라, 같은 답을 지키면서 비용이 어디에서 달라지는지를 설명하는 순서다.

\section*{5. 각 장에서 살펴볼 선택과 비용}

각 장에서는 새로운 장치를 소개하는 데 그치지 않고, 어떤 선택을 했으며 그 선택에 어떤 비용이 따르는지를 다음의 여섯 질문으로 살펴본다.

\begin{center}
\small
\begin{tabular}{>{\raggedright\arraybackslash\sffamily\bfseries}p{0.29\textwidth} >{\raggedright\arraybackslash}p{0.58\textwidth}}
\toprule
보존한 기능·정확한 결과 & \(C=A\times B\)가 뜻하는 수학적 결과 \\
\midrule
계산이 요구하는 양 & 연산 수, 의존 단계, 이동량, 필요한 저장량은 앞으로 각 장에서 하나씩 셈 \\
\midrule
기계가 제공할 자원 & 연산기, 시간, 통로, 가까운 저장공간, 값 하나에 배정한 비트 수는 모두 유한함 \\
\midrule
유한 정밀도에서 달라질 수 있는 것 & 제한된 비트와 계산 순서에 따라 같은 실수식도 조금 다른 기계 결과를 낼 수 있음 \\
\midrule
설계 대가 & 속도·면적·전력·제어·프로그램 가능성을 하나의 비용으로 합치지 않고 각각 따로 비교함 \\
\midrule
아직 해결되지 않은 제약 & 한 줄의 계산을 가장 작은 규칙까지 어떻게 내려가 설명할 것인가 \\
\bottomrule
\end{tabular}
\end{center}

이제 가장 아래의 질문으로 간다. 가능한 입력을 빠짐없이 적었을 때, 우리가 원하는 규칙은 어떻게 하나의 표가 되는가. 제1장은 실제 진리표 한 장에서 시작한다.
